\section{Related Work}
\label{sec:related}

\paragraph*{Tactics in esports}
TAGNet~\cite{zhang2026tagnet} is the closest study and the template we mirror section by section; we replace its learned embedding by objects with units, its fixed $K$ by held-out evidence, and its context-in-the-representation design by a cut.
Spatio-temporal team behaviour in MOBAs was first studied in \emph{Dota~2}, where skill levels differ in how teams spread across zones~\cite{drachen2014skill}.
In \emph{Counter-Strike}, actions have been valued by changes in win probability~\cite{xenopoulos2020valuing} and trajectory data released at scale~\cite{xenopoulos2022esta}.
LoL research has focused on outcome prediction~\cite{hodge2021win,yang2022predicting}, player rating~\cite{maymin2020smart,debois2025pandaskill}, team networks~\cite{mora2019team} and match-level performance indicators~\cite{novak2020performance}.
We use win probability the way action-valuation work does~\cite{decroos2019actions,xenopoulos2020valuing}, but to price the outcomes of a game between teams rather than to score individual actions.

\paragraph*{Space and formations in team sports}
Dominant regions~\cite{taki2000visualization} are Voronoi cells of player positions; we compute them on a navigation graph so that walls count.
Role-based alignment of formations~\cite{lucey2013representing,bialkowski2014large} and the broader spatio-temporal literature~\cite{gudmundsson2017spatio} treat formations as geometric templates.
We treat a formation as a set partition and use the lattice metric~\cite{monjardet1981metrics} to count merges and splits.
Centred partition processes~\cite{paganin2021centered} place priors around a centre partition; HASSE lets the centre move along a path over time and adds rigidity and inertia, inside a sticky HMM~\cite{fox2011sticky}.
Sparse finite mixtures~\cite{fruhwirth2019here} motivate letting the evidence, not a fixed $K$, set the number of patterns.

\paragraph*{Game theory in sport}
Penalty kicks and serves are the canonical tests of equilibrium play by professionals~\cite{walker2001minimax,chiappori2002testing,palacios2003professionals}.
Risk taking by trailing contestants is documented in tournaments~\cite{genakos2012interim} and derived from the convexity of win probability~\cite{cabral2003rd,skinner2011scoring}, the logic behind pulling the goalie~\cite{beaudoin2010pulling}.
Our explanation layer combines a quantal response equilibrium~\cite{mckelvey1995quantal}, whose empirical content requires fixed payoffs~\cite{haile2008empirical}, with Lanchester--Tullock fights~\cite{lanchester1916aircraft,tullock1980efficient} and a Blotto second front~\cite{roberson2006colonel}, estimated as in empirical industrial organisation~\cite{aguirregabiria2007sequential}.
The Shapley decomposition~\cite{shapley1953value} lets the convexity argument compete with the fight law and the second front in one unit.
